\section{拆 L 与 Software 3.0：不是换名词，而是换数据路径}

前面分别讲了软件栈演化和拆 Language，本章把两者放在一起看。刘先明反复提醒，很多新名词本质上做的是类似事情：更大模型、更大数据、更少人工规则、更直接地从输入到动作。所谓 Software 3.0、VLA、VLM、端到端，并不是为了制造概念，而是在寻找更可 scaling 的数据路径。

\subsection{为什么中间结构会反复出现}

本节解释一个工程悖论：中间结构往往是为了解决可控性而加入的，但系统复杂到一定程度后，它们又会反过来限制学习系统的上限。

自动驾驶系统里，中间结构很诱人：感知结果、车道线、目标框、轨迹候选、go point、meta action、language token。它们让系统更可解释、更容易 debug，也更符合传统工程分工。但每增加一个中间结构，就可能增加一个人工设计瓶颈：数据被压缩成某种人为格式，模型只能在这个格式内学习。

\begin{knowledgebox}{中间结构的利弊}
\begin{tabularx}{\textwidth}{p{0.22\textwidth}p{0.34\textwidth}X}
\toprule
维度 & 好处 & 代价 \\
\midrule
可解释性 & 人能看懂模块输出，便于定位问题 & 可能牺牲原始连续信息。 \\
工程分工 & 感知、预测、规控可分团队推进 & 模块接口成为系统上限。 \\
安全上线 & 更容易做规则保护和人工兜底 & 规则越多，数据 scaling 越慢。 \\
数据利用 & 中间标签可以做监督学习 & 人工标签会限制规模和表达力。 \\
\bottomrule
\end{tabularx}
\end{knowledgebox}

\subsection{拆 L 的真正目标：让数据自己说话}

接下来回到小鹏的选择：拆 L 不是为了显得激进，而是为了让更大规模、更原始、更连续的数据进入训练链路。

拆掉 Language Bottleneck 的目标，不是让系统更神秘，而是让传感器、驾驶轨迹和动作反馈能以更少人工变换进入训练。语言可以用于指令和高层语义，但如果每个低层动作都要经过语言 token，它就会把连续世界离散化，阻碍自监督学习和大规模数据利用。

\begin{importantbox}{Software 3.0 的最小定义}
在本讲义中，Software 3.0 指用大模型、大数据、自监督信号和云端模型工厂，把物理世界任务从人工规则系统推进到可持续 scaling 的学习系统。
\end{importantbox}

\subsection{本章小结}

Software 3.0 的关键不是术语，而是数据路径更短。越少人工中间层，越有机会让模型直接从真实数据中学习；但越少中间层，也越需要更强评测、安全和工程治理。

